\section{GDPval: comparar la salida directamente con expertos de la industria}

\subsection{De «poder completarlo» a «si la entrega es suficientemente buena»}

\videofigure{figures/economic-value.jpg}{GDPval no solo pregunta si el agent completó la tarea, sino que hace que expertos comparen los productos de trabajo reales de la IA con los de profesionales pares.}{00:28:06--00:28:38}

\videofigure{figures/gdpval-sectors.jpg}{Las tareas cubren nueve sectores de alto PIB y 44 ocupaciones, en un intento de representar el trabajo de conocimiento común en la economía.}{00:28:38--00:30:02}

\videofigure{figures/gdpval-task-examples.jpg}{Los productos de las tareas incluyen CAD, análisis de inversión, evaluación de imágenes de enfermería, presentaciones, tablas y reportes operativos.}{00:30:02--00:34:16}

A diferencia de los benchmarks de código, las salidas de GDPval a menudo no tienen una respuesta única. La evaluación usa preferencia por pares de expertos de la industria: comparar el producto de la IA con el producto de un experto humano sin conocer el origen, y determinar si gana, pierde o hay empate.

\begin{importantbox}{GDPval mide la calidad del producto, no el proceso de pensamiento oculto}
Los expertos ven el entregable final, así que la métrica se acerca más a la experiencia del cliente; pero no puede indicar directamente si el proceso del agent es confiable, si citó hechos incorrectos o si fallaría con otra entrada.
\end{importantbox}

\subsection{Construcción de los datos y características de las tareas}

\videofigure{figures/gdpval-construction.jpg}{GDPval primero selecciona sectores por PIB, y luego selecciona ocupaciones y actividades de trabajo típicas, con una construcción representativa de arriba hacia abajo.}{00:34:16--00:35:08}

\videofigure{figures/gdpval-task-characteristics.jpg}{El tiempo promedio humano para completar una sola tarea es de aproximadamente siete horas; las tareas son multimodales, dependen de archivos de referencia, y la mayoría no tiene una única respuesta correcta.}{00:35:08--00:36:08}

Este tipo de tareas se parece más al trabajo de oficina que las tareas de software de METR, pero también es más difícil de calificar automáticamente. Una evaluación de alta calidad requiere expertos del dominio, un rubric claro, pruebas a ciegas y verificaciones de consistencia, con un costo mucho mayor que ejecutar pruebas unitarias.

\begin{warningbox}{La evaluación de preferencia se ve afectada por el estilo y la presentación}
Un reporte con buen formato y lenguaje fluido puede ocultar errores de hecho; un producto humano profesional pero sencillo también puede ser subvalorado. Se debe registrar accuracy, instruction following, formatting y severity de manera conjunta, en lugar de mirar solo la preferencia general.
\end{warningbox}

\subsection{El progreso de los modelos no sigue una sola línea}

\videofigure{figures/gdpval-improvement.jpg}{El pairwise win rate en GDPval mejora significativamente entre generaciones de modelos, pero de forma escalonada en lugar de un crecimiento lineal suave.}{00:36:08--00:38:18}

\videofigure{figures/model-specific-strengths.jpg}{Distintos modelos muestran fortalezas claramente diferentes: algunos destacan en estética y documentos, otros son mejores en cálculo, instruction following o herramientas complejas.}{00:38:18--00:39:08}

La clasificación de los modelos depende de la combinación de tareas. Un modelo puede tener ventaja en slides, PDF y disposición visual, mientras que otro puede ser más fuerte en cálculo de tablas, precisión textual o herramientas de código. Por lo tanto, el despliegue empresarial debería enrutar por tarea, en lugar de asumir que un solo modelo domina todas las ocupaciones.

\begin{knowledgebox}{El model routing es un producto directo de la evaluación}
Si el benchmark conserva etiquetas de grano fino sobre ocupación, modalidad y tipo de error, entonces se puede entrenar un router: elegir el modelo, la cadena de herramientas y la intensidad de revisión según la tarea, y no solo usarlo para publicar una tabla general.
\end{knowledgebox}

\subsection{Severidad de los fallos y costo económico}

\videofigure{figures/gdpval-failure-modes.jpg}{Los problemas comunes incluyen errores de instruction following, de formato y de accuracy; la composición de errores difiere entre modelos.}{00:39:08--00:40:08}

\videofigure{figures/failure-severity.jpg}{Muchos fallos son de tipo «aceptable pero no suficientemente bueno»; la proporción de errores realmente catastróficos es baja, pero no se puede ignorar.}{00:40:08--00:41:08}

\videofigure{figures/economic-speed-cost.jpg}{Al considerar «intentar varias veces y luego dejar que un humano corrija», el modelo puede aportar beneficios de velocidad y costo, pero el beneficio depende del costo de reintentar y de revisar.}{00:41:08--00:42:02}

Se puede expresar con un modelo simplificado de costo esperado:

$$
C(n)=nC_{\mathrm{AI}}+
\bigl(1-P_n(\text{acceptable})\bigr)C_{\mathrm{repair}}+C_{\mathrm{review}}.
$$

\begin{itemize}
    \item $n$: número de generaciones o reintentos;
    \item $C_{\mathrm{AI}}$: costo de cada ejecución del modelo;
    \item $P_n$: probabilidad de obtener un resultado aceptable a partir de $n$ intentos;
    \item $C_{\mathrm{repair}}$: costo de que un experto repare un producto fallido;
    \item $C_{\mathrm{review}}$: costo de revisión necesario sin importar si hubo éxito o no.
\end{itemize}

\begin{importantbox}{El beneficio económico proviene del flujo de trabajo «IA + humano»}
Si la revisión y la reparación son económicas, incluso una tasa de éxito baja en un solo intento puede tener valor; si los errores son difíciles de detectar o tienen consecuencias graves, el costo de la revisión humana anulará la ventaja de velocidad del modelo.
\end{importantbox}

\videofigure{figures/performance-variation.jpg}{El desempeño en GDPval varía según el sector, la duración de la tarea y la modalidad; las tareas cortas y algunos sectores se acercan más al nivel de un experto.}{00:42:02--00:45:48}

\subsection{Resumen de la sección}

GDPval hace avanzar la evaluación de «responder correctamente una pregunta» a «entregar un producto profesional comparable». Revela las fortalezas de los modelos, la severidad de los errores y el costo de la colaboración entre humano y máquina, pero aun así se ve afectado por el contexto y las preferencias de los expertos.

